% Generated draft: requires mathematical and source review.
\documentclass[11pt]{article}
\usepackage[margin=25mm]{geometry}
\usepackage{fontspec}
\setmainfont{Latin Modern Roman}
\usepackage{amsmath,amssymb,mathtools,unicode-math}
\setmathfont{Latin Modern Math}
\usepackage{enumitem,needspace,xurl,hyperref}
\setlength{\parindent}{0pt}
\setlength{\parskip}{5pt}
\setlength{\emergencystretch}{3em}
\newcommand{\N}{\mathbb{N}}
\newcommand{\Z}{\mathbb{Z}}
\newcommand{\Q}{\mathbb{Q}}
\newcommand{\R}{\mathbb{R}}
\newcommand{\C}{\mathbb{C}}
\newcommand{\F}{\mathbb{F}}
\newcommand{\A}{\mathbb{A}}
\newcommand{\PP}{\mathbb{P}}
\newcommand{\E}{\mathbb{E}}
\newcommand{\eps}{\varepsilon}
\newcommand{\dd}{\,\mathrm{d}}
\begin{document}
\section*{UnsolvedMath ID: 30003820}
\textbf{Source problem code:} catalogue entry
\subsection{Definitions and mathematical statement}
For a prime power \(q\) and an integer \(n\geq 1\), let
\[
U_n(\F_q)=\{(u_{ij})\in\operatorname{GL}_n(\F_q):u_{ii}=1\text{ for all }i,\ u_{ij}=0\text{ whenever }i>j\}
\]
be the upper unitriangular group. Let \(k(U_n(\F_q))\) denote its number of conjugacy classes.

Let \(T\) and \(x\) be formal indeterminates. Define \(A_n(T)\in\Q[T]\), for \(n\geq0\), by
\[
\sum_{n\geq0}A_n(T)\frac{x^n}{n!}
=(1-\sin(Tx))^{-1/T}
=\exp\!\left(-\frac{\log(1-\sin(Tx))}{T}\right).
\]
The last expression specifies the formal power series interpretation: the numerator is divisible by \(T\), so the apparent division by \(T\) has a removable singularity. In particular, \(A_0(T)=1\).

\textbf{Conditional comparison question.} Fix \(n\geq1\). Suppose that \(S\) is an infinite family of prime powers and that a single polynomial \(P_{n,S}(X)\) satisfies
\[
P_{n,S}(q)=k(U_n(\F_q))\qquad\text{for every }q\in S.
\]
Is it necessarily the case that
\[
P_{n,S}(T+1)-A_n(T)\in\Z_{\geq0}[T]?
\]
Here membership in \(\Z_{\geq0}[T]\) means that every coefficient is a nonnegative integer, including the possibility of zero. The polynomial \(P_{n,S}\), if it exists, is uniquely determined by its values on \(S\). The source notation \(k(U_n(\F_{T+1}))\) is understood as the polynomial substitution \(P_{n,S}(T+1)\), not as a finite field with an indeterminate number of elements.

\textbf{Unresolved source qualifications.} The original statement restricts \(S\) to primes, whereas the clean and displayed statements allow prime powers. These are distinct formulations and neither restriction is selected here as authoritative. The source also leaves unclear whether one family \(S\) is to work simultaneously for every \(n\), whether a separate family may be chosen for each \(n\), and whether the assertion concerns every qualifying family or a particular intended family. The displayed question gives the conditional comparison for a specified \(n\) and \(S\); its full intended quantifier scope requires review. The source does not specify the coefficient ring of the representing polynomial.
\subsection{Short English statement}
The question compares the conjugacy-class counting polynomial of an upper unitriangular group, expressed in \(T=q-1\), with polynomials defined by a formal exponential generating function. It asks whether their difference has nonnegative integer coefficients whenever the counting function agrees with a polynomial on an infinite family. The supplied versions disagree on primes versus prime powers and leave the scope of the family and dimension quantifiers unresolved.
\subsection{Sources}
\begin{itemize}
\item Supplied local UnsolvedMath registry record, permanent ID \(30003820\), title \emph{Positive Polynomial Comparisons for Unitriangular Conjugacy Classes}, exact source problem code \texttt{catalogue entry}, in the Oberwolfach Reports Open Problems collection. No individual proposer is identified; the supplied proposed year is \(2018\).
\item Supplied source citation: \emph{Enumerative Combinatorics (2019)}, Oberwolfach Reports, pp.\ 1381--1464, DOI \texttt{10.4171/owr/2018/23}. Original OWR report: \url{https://doi.org/10.4171/owr/2018/23}. The supplied citation uses \(2019\), while the DOI and year metadata refer to \(2018\); their relationship has not been independently checked.
\item Supplied literature reference, labeled ``Pak--Soffer on Higman's conjecture'': \url{https://arxiv.org/abs/1507.00411}.
\item Literature assessment checked at \texttt{2026-08-21}: this is supplied dated curation, not an independently verified assessment. It labels the problem open, reports verification of the proposed comparison only through \(n=16\), and states that the underlying Higman polynomiality question remains unresolved and that no proof or counterexample for this stronger comparison was found in the literature considered. The supplied background expressly requests independent verification of time-sensitive status and cited claims.
\item The registry labels the statement \texttt{corrected\_verified} and describes its editing as repairs to notation, typography, and sentence structure preserving mathematical content. This supplied editorial label does not establish mathematical correctness or resolve the formulation discrepancies above.
\end{itemize}
\end{document}
